\documentclass[12pt,a4paper]{article}
\usepackage{amsmath,amssymb,amsthm}
\usepackage{hyperref}
\usepackage{geometry}
\geometry{margin=2.5cm}
\usepackage{enumitem}
\usepackage{booktabs}

\newtheorem{theorem}{Theorem}[section]
\newtheorem{lemma}[theorem]{Lemma}
\newtheorem{corollary}[theorem]{Corollary}
\newtheorem{proposition}[theorem]{Proposition}
\theoremstyle{definition}
\newtheorem{definition}[theorem]{Definition}
\theoremstyle{remark}
\newtheorem{remark}[theorem]{Remark}

\title{Ledger Realism:\\
The Recognition Science Interpretation of Quantum Mechanics}

\author{Jonathan Washburn\\
Recognition Science Research Institute\\
Austin, Texas\\
\texttt{washburn.jonathan@gmail.com}}

\date{March 2026}

\begin{document}
\maketitle

\begin{abstract}
We present \emph{Ledger Realism}, the Recognition Science (RS)
interpretation of quantum mechanics.  The fundamental ontology is
the recognition ledger---a discrete, finite, deterministic
substrate that records every recognition event.  The past
(committed ledger entries) is real and definite.  The present
(current tick) is privileged.  The future (uncommitted entries) is
genuinely open.  The wave function $\psi$ is the amplitude for the
next ledger entry---epistemic at the inter-tick level, but the
ledger itself is ontic.  Superposition describes inter-tick
evolution between definite ledger states, not simultaneously real
alternatives.  This interpretation is neither Copenhagen (no
observer dependence), nor Many-Worlds (no branching), nor Bohmian
(no pilot wave).  It satisfies Bell's desiderata for a ``beable''
theory: the local beables are ledger entries, discrete and
localized in spacetime.  The interpretation resolves the
measurement problem (collapse = commit), the basis problem
(pointer states selected by $J$-cost minimization), and the
preferred-basis problem (einselection from environmental $J$-cost
coupling).
\end{abstract}

%% ============================================================
\section{Introduction}
\label{sec:intro}
%% ============================================================

The interpretation of quantum mechanics remains one of the deepest
questions in physics~\cite{Bell1987}.  Every interpretation must
address four questions:
\begin{enumerate}[label=(\roman*)]
  \item What is the ontology? (What is ``real''?)
  \item What is the status of $\psi$?
  \item When and how does ``collapse'' occur?
  \item What determines definite outcomes?
\end{enumerate}
Ledger Realism, grounded in the RS forcing chain~\cite{RS_Axioms},
provides precise answers to all four.

%% ============================================================
\section{The Ontology}
\label{sec:ontology}
%% ============================================================

\begin{theorem}[Ledger Ontology]
\label{thm:ontology}
The recognition ledger is the fundamental reality.  It has three
defining properties:
\begin{enumerate}[label=(\roman*)]
  \item \textbf{Discrete}: composed of voxels (size $\ell_0$)
  updated at ticks (duration $\tau_0$).
  \item \textbf{Finite}: bounded by the Hubble volume
  ($\sim 10^{183}$ voxels).
  \item \textbf{Definite at each tick}: each voxel contains
  exactly one integer-valued entry per tick.
\end{enumerate}
\end{theorem}

\begin{remark}
The ledger is not a simulation; it is not running ``on'' anything
else.  It is the fundamental substrate.  There is no deeper level.
\end{remark}

%% ============================================================
\section{Temporal Structure}
\label{sec:time}
%% ============================================================

\begin{theorem}[A-Theory of Time]
\label{thm:time}
\begin{itemize}
  \item \textbf{Past}: committed ledger entries---real, definite,
  irreversible.  The past is a growing chain of frozen records.
  \item \textbf{Present}: the current tick---privileged, the
  locus of recognition.  Only the present tick is ``active.''
  \item \textbf{Future}: uncommitted entries---genuinely open,
  not yet determined, not hidden.
\end{itemize}
\end{theorem}

\begin{remark}
This is an A-theory of time (past--present--future are
ontologically distinct), in contrast to the B-theory (block
universe) implicit in most interpretations of quantum mechanics.
The block universe is incompatible with RS because uncommitted
future entries do not yet exist as ledger records.
\end{remark}

%% ============================================================
\section{The Status of $\psi$}
\label{sec:psi}
%% ============================================================

\begin{theorem}[Wave Function as Computational Amplitude]
\label{thm:psi}
$\psi$ is the amplitude for the next ledger entry.  It is:
\begin{enumerate}[label=(\roman*)]
  \item \textbf{Epistemic} at the inter-tick level: it encodes
  the probabilities of upcoming commits.
  \item \textbf{Not ontic}: $\psi$ is not an element of physical
  reality.  The ledger entries are physical; $\psi$ is
  computational.
  \item \textbf{Predictive}: $P(i) = |\psi_i|^2$ (Born rule).
  \item \textbf{Complex}: $\psi \in \mathbb{C}$ because the
  8-tick phase structure requires complex amplitudes.
\end{enumerate}
\end{theorem}

\begin{remark}
The relationship between $\psi$ and the ledger is analogous to
the relationship between a probability distribution and the
outcomes it predicts.  The distribution is useful but not
``real''; the outcomes are real.
\end{remark}

%% ============================================================
\section{Collapse = Commit}
\label{sec:collapse}
%% ============================================================

\begin{definition}[Ledger Commit]
At each 8-tick epoch boundary ($t \bmod 8 = 0$), the ledger
transitions from the uncommitted (superposition) phase to the
committed (definite) phase.  One branch is actualized; the others
are discarded.
\end{definition}

\begin{theorem}[Collapse Timing]
``Collapse'' occurs at every 8-tick epoch boundary---it is
automatic, regular, and universal.  No observer, no consciousness,
no macroscopic apparatus is required.
\end{theorem}

\begin{theorem}[Irreversibility of Commit]
A committed ledger entry cannot be uncommitted.  This is because
reversal would create an unbalanced double-entry (credit without
debit), which has infinite $J$-cost.
\end{theorem}

%% ============================================================
\section{Bell's Beables}
\label{sec:beables}
%% ============================================================

\begin{theorem}[Ledger Entries as Local Beables]
Bell argued that a satisfactory interpretation must specify
\emph{local beables}---elements of physical reality localized in
spacetime~\cite{Bell1987}.  In Ledger Realism, the local beables
are the committed ledger entries: discrete, localized (at specific
voxels and ticks), and definite (integer-valued).
\end{theorem}

\begin{remark}
Copenhagen has no beables (ontology is unclear).  Many-Worlds
has $\psi$ as a non-local beable.  Bohmian mechanics has particle
positions.  Ledger Realism has ledger entries---the most concrete
and discrete of all proposed beables.
\end{remark}

%% ============================================================
\section{The Born Rule}
\label{sec:born}
%% ============================================================

\begin{theorem}[Born Rule from $J$-Cost Phase Invariance]
\begin{equation}
  \boxed{P(i) = |\psi_i|^2.}
\end{equation}
\end{theorem}

\begin{proof}
The cost functional $J(x) = \tfrac{1}{2}(x + x^{-1}) - 1$ depends
only on $|x|$, not on the phase: $J(re^{i\theta}) = J(r)$.  The
probability of a branch must be non-negative, normalized, and
phase-independent.  The unique such function of the amplitude is
$P(i) = |c_i|^2$.
\end{proof}

%% ============================================================
\section{Comparison with Other Interpretations}
\label{sec:comparison}
%% ============================================================

\begin{center}
\renewcommand{\arraystretch}{1.3}
\begin{tabular}{@{}lccccc@{}}
\toprule
& \textbf{Copen.} & \textbf{MWI} & \textbf{Bohm} &
\textbf{GRW} & \textbf{RS} \\
\midrule
Ontology & Unclear & $\psi$ & Part.\ + $\psi$ & $\psi$ +
collapse & Ledger \\
$\psi$ status & Epist. & Ontic & Ontic & Ontic & Epist. \\
Collapse & Observer & None & Effective & Spontaneous &
Commit \\
Determinism & No & Yes & Yes & No & Tick-level \\
Time & B-theory & B-theory & Either & B-theory & A-theory \\
Beables & None & $\psi$ & Positions & $\psi$ & Entries \\
Free params & Many & 0 & 0 & 2 ($\lambda$, $r_c$) & 0* \\
\bottomrule
\end{tabular}
\end{center}

\noindent *The core RS ontology (ledger, Born rule, collapse mechanism)
has zero free parameters; all are fixed by $\varphi$.  The
environmental coupling $g$ entering the decoherence timescale
$\tau_{\mathrm{dec}} = \tau_0 \varphi^{-Ng}$ is expected to be
derivable from the RS Hamiltonian for each environment type; this
derivation is an identified open problem.

%% ============================================================
\section{Addressing Objections}
\label{sec:objections}
%% ============================================================

\begin{remark}[``Is this just Copenhagen?'']
No.  Copenhagen makes collapse observer-dependent and leaves the
ontology unspecified.  Ledger Realism specifies the ontology
(ledger entries) and makes collapse automatic (every 8-tick
epoch).
\end{remark}

\begin{remark}[``Is this a hidden-variable theory?'']
Not in the Bell sense.  The ledger entries are not ``hidden''---they
\emph{are} the physical reality.  The wave function is the hidden
(epistemic) quantity; the entries are the manifest (ontic) ones.
\end{remark}

\begin{remark}[``Does this violate Bell's theorem?'']
No.  Bell's theorem shows that no \emph{local} hidden-variable
theory can reproduce quantum correlations.  The ledger is
\emph{global}: entangled particles share ledger entries (non-local
constraints).  This produces quantum correlations without
superluminal signaling.
\end{remark}

%% ============================================================
\section{Predictions}
\label{sec:predictions}
%% ============================================================

\begin{enumerate}
  \item \textbf{No collapse modifications}: GRW-type spontaneous
  collapse models (with parameters $\lambda$ and $r_c$) are
  excluded.  Collapse is deterministic (8-tick epoch), not
  stochastic.

  \item \textbf{Born rule exact}: No deviations from
  $P = |\psi|^2$ at any scale.

  \item \textbf{No branching}: No evidence for Everettian
  many-worlds branching will be found.

  \item \textbf{Decoherence timescale}: Macroscopic objects
  ($N \sim 10^{23}$ modes, $g \sim 0.1$) decohere in
  $\tau_{\mathrm{dec}} = \tau_0 \cdot \varphi^{-Ng}$, with
  $Ng \sim 10^{22}$---an exponent so large the time is
  effectively zero (see companion paper on quantum-to-classical
  transition for the precise formula and $g$ discussion).
\end{enumerate}

%% ============================================================
\section{Conclusion}
\label{sec:conclusion}
%% ============================================================

Ledger Realism provides a complete, coherent interpretation of
quantum mechanics: the ledger is real; the wave function is
computational.  The past is definite, the present is privileged,
the future is open.  Collapse is the automatic commit operation.
Bell's beables are ledger entries.  The Born rule follows from
$J$-cost phase invariance.  The core theory has zero free
parameters (all fixed by $\varphi$); only the environmental
coupling $g$ remains to be derived from the RS Hamiltonian for
specific interaction types.

%% ============================================================
\begin{thebibliography}{99}

\bibitem{RS_Axioms}
J.~Washburn, ``Recognition Science: Foundations,''
RS Axioms Document (2026).

\bibitem{Bell1987}
J.~S.~Bell, \textit{Speakable and Unspeakable in Quantum
Mechanics} (Cambridge University Press, 1987).

\end{thebibliography}

\end{document}
